\documentclass[a4paper,10.5pt]{article}
\usepackage[T1]{fontenc}\usepackage[utf8]{inputenc}
\usepackage[top=2.2cm,bottom=2.2cm,left=2.35cm,right=2.35cm]{geometry}
\usepackage{hyperref}\usepackage{xcolor}\definecolor{accent}{HTML}{1F4E78}\hypersetup{colorlinks=true,urlcolor=accent}
\pagestyle{empty}\setlength{\parindent}{0pt}\setlength{\parskip}{7pt}
\begin{document}
\textbf{Prabhavit Gupta}\\
Venice, Italy \textbar{} +39 351 610 7807 \textbar{} prabhavitg@gmail.com\\[10pt]
Dear Wiz-Team hiring team,

I am applying for the Technical Support Engineer role in Turin. I have three years of professional experience as a full-stack and AI developer, most of it spent close to production incidents, API integrations and the kind of troubleshooting this role needs.

At FlowSpace, a live workspace marketplace, I built support agents that answer customer queries directly from live platform data and escalate to a human when confidence is low, so I have hands on experience designing for graceful failure and knowing when to hand off. I also shipped REST integrations with OfficeRnD and Nexudus for automated bookings, which means reading logs, API requests and third-party responses to trace an issue back to its root cause is part of my regular work, not a new skill.

At Clubee, a SaaS platform in Luxembourg, I integrated the Clubee API across the frontend and ran unit testing and QA on every new feature, which built the habit of checking application behaviour carefully before and after a change ships, the same discipline production support demands.

I work daily with Python and TypeScript, am comfortable with SQL for corrective queries, and use tools like Jira and GitHub as part of normal delivery. I am based in Venice, authorised to work in Italy without sponsorship, and open to the Turin hybrid schedule and the shifted hours this role covers.

Regards,\\
Prabhavit Gupta
\end{document}
